\subsection{Measures, Charges, and Basic Constructions}

\begin{definition} \label{def:measure-measurability-core-source-todo}\label{def:ring-sigma-ring-set-function} \hlt{Rings, $\sigma$-Rings, and Set Functions}\\
Let $\mathcal{A}\subseteq\mathcal{P}(X)$.
\begin{enumerate}[label=(\roman*)]
\setlength{\itemsep}{0pt}
\item $\mathcal{A}$ is a \hlt{ring of sets} iff $\emptyset\in\mathcal{A}$ and $A,B\in\mathcal{A}\Rightarrow A\cup B,\ A\setminus B\in\mathcal{A}$
\item $\mathcal{A}$ is a \hlt{$\sigma$-ring} iff it is a ring and $(A_n)\subseteq\mathcal{A}\Rightarrow\bigcup_{n=1}^{\infty}A_n\in\mathcal{A}$
\item A \hlt{set function} on $\mathcal{A}$ is $\phi:\mathcal{A}\to[-\infty,+\infty]$ with $\{+\infty,-\infty\}\nsubseteq\phi(\mathcal{A})$.
\item $\phi$ is \hlt{additive} iff $A\cap B=\emptyset\Rightarrow\phi(A\cup B)=\phi(A)+\phi(B)$
\item $\phi$ is \hlt{countably additive} iff $(A_n)\subseteq\mathcal{A}$ pairwise disjoint, $\bigcup_nA_n\in\mathcal{A}$ $\Rightarrow$ $\phi\left(\bigcup_{n=1}^{\infty}A_n\right)=\sum_{n=1}^{\infty}\phi(A_n)$.
\end{enumerate}
\hlt{Structural Commentary}: Rings support finite operations; $\sigma$-rings support countable unions.
\end{definition}

\begin{definition} \label{def:abstract-measure-on-measurable-space} \hlt{Measures on Measurable Spaces}\\
Let $(X,\mathcal{M})$ be a measurable space. A \hlt{measure} on $(X,\mathcal{M})$ is a map $\mu:\mathcal{M}\to[0,+\infty]$ such that:
\begin{enumerate}[label=(\roman*)]
\setlength{\itemsep}{0pt}
\item $\mu(\emptyset)=0$;
\item if $(E_n)\subseteq\mathcal{M}$ is pairwise disjoint, then $\mu\left(\bigcup_{n=1}^{\infty}E_n\right)=\sum_{n=1}^{\infty}\mu(E_n)$
\end{enumerate}
The triple $(X,\mathcal{M},\mu)$ is a \hlt{measure space}.\\
\hlt{Key Insight}: Measure is countably additive size on the chosen observable sets.
\end{definition}

\begin{definition} \label{def:abstract-measures-finite-sigma-finite} \hlt{Finite and $\sigma$-Finite Measures}\\
Let $(X,\mathcal{M},\mu)$ be a measure space.
\begin{enumerate}[label=(\roman*)]
\setlength{\itemsep}{0pt}
\item For nonnegative terms, $\sum_{n=1}^{\infty}a_n\defeq\lim_{N\to\infty}\sum_{n=1}^{N}a_n\in[0,+\infty]$.
\item $\mu$ is \hlt{finite} iff $\mu(X)<+\infty$.
\item $\mu$ is \hlt{$\sigma$-finite} iff there exist $E_n\in\mathcal{M}$ such that $X=\bigcup_{n=1}^{\infty}E_n$ and $\mu(E_n)<+\infty$ for all $n$.
\item A measurable set $E\in\mathcal{M}$ has \hlt{finite measure} iff $\mu(E)<+\infty$.
\item A measurable set $E\in\mathcal{M}$ is \hlt{$\sigma$-finite} iff $E=\bigcup_{n=1}^{\infty}E_n$ for some $E_n\in\mathcal{M}$ with $\mu(E_n)<+\infty$ for all $n$.
\end{enumerate}
\hlt{Key Insight}: $\sigma$-finiteness replaces one finite global bound by a countable finite-measure exhaustion.
\end{definition}

\begin{example} \label{ex:first-abstract-measures} \hlt{First Abstract Measures}
\begin{enumerate}[label=(\roman*)]
\setlength{\itemsep}{0pt}
\item The zero set function $\mu(E)=0$ on $\mathcal{M}$ is a finite measure.
\item If $X\neq\emptyset$ and $\mathcal{M}=2^X$, $\mu(\emptyset)=0$ and $\mu(E)=+\infty$ for $E\neq\emptyset$ give a measure, neither finite nor $\sigma$-finite.
\item For $p\in X$, the \hlt{unit mass} $\delta_p(E)\defeq\mathbf{1}_E(p)$ on $2^X$ is a finite measure.
\item On $2^\N$, \hlt{counting measure} is $\#(E)$ for finite $E$ and $+\infty$ otherwise; it is $\sigma$-finite.\\
More generally, $(a_n)\subseteq[0,+\infty]$ gives $\mu(E)\defeq\sum_{n\in E}a_n$, and every measure on $2^\N$ is uniquely of this form.
\item If $X$ is uncountable and $\mathcal{M}=2^X$, $\mu(E)=0$ for countable $E$ and $\mu(E)=+\infty$ otherwise defines a measure.
\item If $X$ is uncountable and $\mathcal{C}_X=\{E\subseteq X:E\text{ is countable or }X\setminus E\text{ is countable}\}$, then
$\mu(E)=0$ on countable sets and $\mu(E)=1$ on cocountable sets defines a measure on $\mathcal{C}_X$.
\item On $2^\N$, $\mu(E)=0$ for finite $E$ and $\mu(E)=+\infty$ otherwise is not a measure, since $\N=\bigcup_{n=1}^{\infty}\{n\}$ gives $+\infty\neq0$; the constant $+\infty$ rule fails at $\emptyset$.
\item Lebesgue measure on $\mathcal{B}(\R)$ is $\sigma$-finite and gives length $b-a$ to standard intervals; Borel--Stieltjes measures come from increasing right-continuous functions as in Theorem \ref{thm:borel-stieltjes-measure-generation}.
\end{enumerate}
\hlt{Structural Commentary}: Measures may be zero, atomic, counting-type, cardinality-based, or geometric.\\
\hlt{Key Insight}: Countable additivity turns disjoint decompositions into nonnegative series.
\end{example}

\begin{proposition} \label{prop:elementary-measure-constructions} \hlt{Elementary Measure Constructions}\\
Let $(X,\mathcal{M})$ be a measurable space.
\begin{enumerate}[label=(\roman*)]
\setlength{\itemsep}{0pt}
\item If $\mu$ is a measure on $\mathcal{M}$ and $A\in\mathcal{M}$, then $\mu_A(E)\defeq\mu(A\cap E)$ defines a measure $\mu_A$ on $\mathcal{M}$.
\item If $\mu_1,\ldots,\mu_m$ are measures on $\mathcal{M}$ and $a_1,\ldots,a_m\geq0$, then $\lambda\defeq\sum_{j=1}^{m}a_j\mu_j$ is a measure on $\mathcal{M}$.
\item If $(\mu_n)_{n\geq1}$ are measures on $\mathcal{M}$ with $\mu_n(X)=1$, then $\lambda\defeq\sum_{n=1}^{\infty}2^{-n}\mu_n$ is a measure on $\mathcal{M}$ and $\lambda(X)=1$.
\end{enumerate}
\begin{proof}
\begin{enumerate}[label=(\roman*)]
\setlength{\itemsep}{0pt}
\item $\mu_A(\emptyset)=0$. If $(E_k)$ is disjoint in $\mathcal{M}$, then $A\cap\bigsqcup_kE_k=\bigsqcup_k(A\cap E_k)$, whence $\mu_A(\bigsqcup_kE_k)=\sum_k\mu_A(E_k)$.
\item $\lambda(\emptyset)=0$. For disjoint $(E_k)$, finite summation gives $\lambda(\bigsqcup_kE_k)=\sum_{j=1}^{m}a_j\sum_k\mu_j(E_k)=\sum_k\lambda(E_k)$.
\item $\lambda(\emptyset)=0$. For disjoint $(E_k)$, nonnegative reordering gives $\lambda(\bigsqcup_kE_k)=\sum_n2^{-n}\sum_k\mu_n(E_k)=\sum_k\lambda(E_k)$.\\
Also $\lambda(X)=\sum_n2^{-n}\mu_n(X)=\sum_n2^{-n}=1$.
\end{enumerate}
\hlt{Structural Commentary}: New measures come from restriction, superposition, and countable convex mixtures.\\
\hlt{Key Insight}: Nonnegative series make the bookkeeping order-independent.
\end{proof}
\end{proposition}

\begin{proposition} \label{prop:measure-basic-laws} \hlt{Basic Laws for Measures}\\
Let $(X,\mathcal{M},\mu)$ be a measure space.
\begin{enumerate}[label=(\roman*)]
\setlength{\itemsep}{0pt}
\item For every $n\geq1$, if $E_1,\ldots,E_n\in\mathcal{M}$ are pairwise disjoint, then $\mu\left(\bigcup\limits_{k=1}^{n}E_k\right)=\sum\limits_{k=1}^{n}\mu(E_k)$.
\item If $A,B\in\mathcal{M}$ and $A\subseteq B$, then $\mu(A)\leq\mu(B)$.
\item If $A,B\in\mathcal{M}$, $A\subseteq B$, and $\mu(A)<+\infty$, then $\mu(B\setminus A)=\mu(B)-\mu(A)$.\\
In particular, if $\mu(A)=0$, then $\mu(B\setminus A)=\mu(B)$.
\item If $E\in\mathcal{M}$ and $(E_k)_{k\geq1}\subseteq\mathcal{M}$ satisfy $E\subseteq\bigcup\limits_{k=1}^{\infty}E_k$, then $\mu(E)\leq\sum\limits_{k=1}^{\infty}\mu(E_k)$.
\end{enumerate}
\begin{proof}
\begin{enumerate}[label=(\roman*)]
\setlength{\itemsep}{0pt}
\item Define $G_k=E_k$ for $1\leq k\leq n$ and $G_k=\emptyset$ for $k>n$. Then $(G_k)$ is pairwise disjoint and countable additivity gives $\mu\left(\bigcup\limits_{k=1}^{n}E_k\right)=\mu\left(\bigcup\limits_{k=1}^{\infty}G_k\right)=\sum\limits_{k=1}^{\infty}\mu(G_k)=\sum\limits_{k=1}^{n}\mu(E_k)$.
\item Since $B=A\sqcup(B\setminus A)$, (i) gives $\mu(B)=\mu(A)+\mu(B\setminus A)$, and $\mu(B\setminus A)\geq0$, hence $\mu(A)\leq\mu(B)$.
\item Since $B=A\sqcup(B\setminus A)$, (i) gives $\mu(B)=\mu(A)+\mu(B\setminus A)$. Because $\mu(A)<+\infty$, subtraction yields $\mu(B\setminus A)=\mu(B)-\mu(A)$. If $\mu(A)=0$, this becomes $\mu(B\setminus A)=\mu(B)$.
\item Set $F_1=E\cap E_1$ and $F_k=E\cap\left(E_k\setminus\bigcup\limits_{j=1}^{k-1}E_j\right)$ for $k\geq2$. Then $(F_k)_{k\geq1}\subseteq\mathcal{M}$ is pairwise disjoint, $\bigcup\limits_{k=1}^{\infty}F_k=E$, and $F_k\subseteq E_k$ for all $k$. Countable additivity and (ii) give $\mu(E)=\sum\limits_{k=1}^{\infty}\mu(F_k)\leq\sum\limits_{k=1}^{\infty}\mu(E_k)$.
\end{enumerate}
\hlt{Structural Commentary}: Finite additivity yields order and excision; disjointization turns covers into sums.\\
\hlt{Key Insight}: Measure disjoint pieces, then compare.
\end{proof}
\end{proposition}

\begin{theorem} \label{thm:continuity-from-below-set-functions} \hlt{Continuity from Below}\\
Let $\mathcal{A}$ be a ring of subsets of $X$, and let $\phi$ be a countably additive set function on $\mathcal{A}$.\\
If $(A_n)_{n\geq1}\subseteq\mathcal{A}$ and $A\in\mathcal{A}$ satisfy $A_n\subseteq A_{n+1}$ for all $n$ and $A=\bigcup_{n=1}^{\infty}A_n$, then $\lim_{N\to\infty}\phi(A_N)=\phi(A)$.
\begin{proof} Set $B_1=A_1$ and $B_n=A_n\setminus A_{n-1}$ for $n\geq2$. Since $\mathcal{A}$ is a ring, each $B_n\in\mathcal{A}$.\\
The sets $(B_n)$ are pairwise disjoint, $A_N=\bigsqcup_{n=1}^{N}B_n$, and $A=\bigsqcup_{n=1}^{\infty}B_n$.\\
Countable additivity, applied to the finite union by appending empty sets, gives $\phi(A_N)=\sum_{n=1}^{N}\phi(B_n)$, while countable additivity gives $\phi(A)=\sum_{n=1}^{\infty}\phi(B_n)$. Thus $\lim_{N\to\infty}\phi(A_N)=\phi(A)$.\\
\hlt{Structural Commentary}: Increasing chains split into disjoint increments.\\
\hlt{Key Insight}: Continuity is convergence of partial sums.
\end{proof}
\end{theorem}

\begin{definition} \label{def:charges-finite-signed-measures} \hlt{Charges}\\
Let $(X,\mathcal{M})$ be a measurable space. A \hlt{charge} is a map $\lambda:\mathcal{M}\to\R$ such that:
\begin{enumerate}[label=(\roman*)]
\setlength{\itemsep}{0pt}
\item $\lambda(\emptyset)=0$.
\item For every pairwise disjoint $(E_n)\subseteq\mathcal{M}$, $\sum\limits_{n=1}^{\infty}\lambda(E_n)$ converges in $\R$ and $\lambda\left(\bigcup\limits_{n=1}^{\infty}E_n\right)=\sum\limits_{n=1}^{\infty}\lambda(E_n)$.
\end{enumerate}
\hlt{Key Insight}: A charge is a real-valued countably additive set function.
\end{definition}

\begin{definition} \label{def:charges-and-variation-measures} \hlt{Variation Candidates of a Charge}\\
Let $\lambda$ be a charge on $\mathcal{M}$. For $E\in\mathcal{M}$, define
\begin{enumerate}[label=(\roman*)]
\setlength{\itemsep}{0pt}
\item $\mathfrak{P}_{\mathcal{M}}(E)\defeq\{(A_j)_{j=1}^{n}:n\in\N,\ E=\bigsqcup_{j=1}^{n}A_j,\ A_j\in\mathcal{M}\}$, the finite measurable partitions of $E$;
\item $\pi_\lambda(E)\defeq\sup\{\lambda(A):A\subseteq E,\ A\in\mathcal{M}\}$, the \hlt{positive variation} candidate;
\item $\nu_\lambda(E)\defeq\sup\{\sum_{j=1}^{n}\abs{\lambda(A_j)}:(A_j)_{j=1}^{n}\in\mathfrak{P}_{\mathcal{M}}(E)\}$, the \hlt{total variation} candidate.
\end{enumerate}
For $E=\emptyset$, the empty partition is allowed and contributes $0$.
\end{definition}

\begin{proposition} \label{prop:charges-linear-combinations-and-variation} \hlt{Charge Operations and Variation Measures}\\
Let $(X,\mathcal{M})$ be a measurable space, let $m\geq1$, and let $\lambda,\lambda_1,\ldots,\lambda_m$ be charges on $\mathcal{M}$.
\begin{enumerate}[label=(\roman*)]
\setlength{\itemsep}{0pt}
\item If $a_1,\ldots,a_m\in\R$ and $\Lambda:\mathcal{M}\to\R$ is defined by $\Lambda(E)\defeq\sum\limits_{j=1}^{m}a_j\lambda_j(E)$, then $\Lambda$ is a charge.
\item If $\mu$ is a measure on $\mathcal{M}$, $f\in\mathscr{L}^1(\mu)$, and $\lambda_f:\mathcal{M}\to\R$ is defined by $\lambda_f(A)\defeq\int_A f\,d\mu$, then $\lambda_f$ is a charge.
\item $\pi_\lambda:\mathcal{M}\to[0,+\infty]$ is a measure.
\item $\nu_\lambda:\mathcal{M}\to[0,+\infty]$ is a measure.
\end{enumerate}
\begin{proof}
\begin{enumerate}[label=(\roman*)]
\setlength{\itemsep}{0pt}
\item $\Lambda(\emptyset)=0$. For $E=\bigsqcup\limits_nE_n$, $\sum\limits_n\Lambda(E_n)$ converges as $\sum\limits_{j=1}^{m}a_j\sum\limits_n\lambda_j(E_n)$; $\Lambda(E)=\sum\limits_{j=1}^{m}a_j\sum\limits_n\lambda_j(E_n)=\sum\limits_n\Lambda(E_n)$.
\item $\lambda_f(\emptyset)=0$. For disjoint $(A_n)$, $\sum\limits_n\abs{\lambda_f(A_n)}\leq\int_{\bigsqcup\limits_nA_n}\abs{f}\,d\mu<+\infty$\\
and $\lambda_f(\bigsqcup\limits_nA_n)=\int_{\bigsqcup\limits_nA_n}f\,d\mu=\sum\limits_n\int_{A_n}f\,d\mu=\sum\limits_n\lambda_f(A_n)$.
\item $\pi_\lambda(\emptyset)=0$ and $\pi_\lambda\geq0$. Let $E=\bigsqcup\limits_nE_n$. For measurable $A\subseteq E$, $\lambda(A)=\sum\limits_n\lambda(A\cap E_n)\leq\sum\limits_n\pi_\lambda(E_n)$, so $\pi_\lambda(E)\leq\sum\limits_n\pi_\lambda(E_n)$. Put $S_N=\sum\limits_{n=1}^{N}\pi_\lambda(E_n)$. If $S_N<+\infty$, choose measurable $A_n\subseteq E_n$ for $n\leq N$ with $\lambda(A_n)>\pi_\lambda(E_n)-\frac{\varepsilon}{N}$; then $\pi_\lambda(E)>S_N-\varepsilon$. If $S_N=+\infty$, some $E_\ell$ has $\pi_\lambda(E_\ell)=+\infty$, hence $\pi_\lambda(E)=+\infty$. Thus $\pi_\lambda(E)\geq S_N$ for all $N$, and $N\to\infty$ gives equality.
\item $\nu_\lambda(\emptyset)=0$ and $\nu_\lambda\geq0$. Let $E=\bigsqcup\limits_nE_n$ and $v(P)=\sum\limits_{A\in P}\abs{\lambda(A)}$. Refining $P=(A_j)_{j=1}^{r}\in\mathfrak{P}_{\mathcal{M}}(E)$ by $E_n$ gives $v(P)\leq\sum\limits_{j=1}^{r}\sum\limits_n\abs{\lambda(A_j\cap E_n)}\leq\sum\limits_n\nu_\lambda(E_n)$, so $\nu_\lambda(E)\leq\sum\limits_n\nu_\lambda(E_n)$. Put $T_N=\sum\limits_{n=1}^{N}\nu_\lambda(E_n)$.\\
If $T_N<+\infty$, choose $P_n\in\mathfrak{P}_{\mathcal{M}}(E_n)$ for $n\leq N$ with $v(P_n)>\nu_\lambda(E_n)-\frac{\varepsilon}{N}$.\\
Adjoining the tail $\bigcup\limits_{n>N}E_n$ gives $\nu_\lambda(E)>T_N-\varepsilon$. If $T_N=+\infty$, some $E_\ell$ has $\nu_\lambda(E_\ell)=+\infty$, and partitions of $E_\ell$ extend to $E$, so $\nu_\lambda(E)=+\infty$. Thus $\nu_\lambda(E)\geq T_N$ for all $N$, and $N\to\infty$ gives equality.
\end{enumerate}
\hlt{Structural Commentary}: Suprema preserve additivity after refinement.\\
\hlt{Key Insight}: Refine first, then take suprema.
\end{proof}
\end{proposition}
